% Options for packages loaded elsewhere
\PassOptionsToPackage{unicode}{hyperref}
\PassOptionsToPackage{hyphens}{url}
\PassOptionsToPackage{dvipsnames,svgnames,x11names}{xcolor}
%
\documentclass[
  11pt,
  a4paper,
]{ltjsarticle}
\usepackage{amsmath,amssymb}
\usepackage{iftex}
\ifPDFTeX
  \usepackage[T1]{fontenc}
  \usepackage[utf8]{inputenc}
  \usepackage{textcomp} % provide euro and other symbols
\else % if luatex or xetex
  \usepackage{unicode-math} % this also loads fontspec
  \defaultfontfeatures{Scale=MatchLowercase}
  \defaultfontfeatures[\rmfamily]{Ligatures=TeX,Scale=1}
\fi
\usepackage{lmodern}
\ifPDFTeX\else
  % xetex/luatex font selection
\fi
% Use upquote if available, for straight quotes in verbatim environments
\IfFileExists{upquote.sty}{\usepackage{upquote}}{}
\IfFileExists{microtype.sty}{% use microtype if available
  \usepackage[]{microtype}
  \UseMicrotypeSet[protrusion]{basicmath} % disable protrusion for tt fonts
}{}
\makeatletter
\@ifundefined{KOMAClassName}{% if non-KOMA class
  \IfFileExists{parskip.sty}{%
    \usepackage{parskip}
  }{% else
    \setlength{\parindent}{0pt}
    \setlength{\parskip}{6pt plus 2pt minus 1pt}}
}{% if KOMA class
  \KOMAoptions{parskip=half}}
\makeatother
\usepackage{xcolor}
\usepackage[margin=24mm]{geometry}
\setlength{\emergencystretch}{3em} % prevent overfull lines
\providecommand{\tightlist}{%
  \setlength{\itemsep}{0pt}\setlength{\parskip}{0pt}}
\setcounter{secnumdepth}{5}
\ifLuaTeX
  \usepackage{selnolig}  % disable illegal ligatures
\fi
\IfFileExists{bookmark.sty}{\usepackage{bookmark}}{\usepackage{hyperref}}
\IfFileExists{xurl.sty}{\usepackage{xurl}}{} % add URL line breaks if available
\urlstyle{same}
\hypersetup{
  colorlinks=true,
  linkcolor={blue},
  filecolor={Maroon},
  citecolor={Blue},
  urlcolor={blue},
  pdfcreator={LaTeX via pandoc}}

\author{}
\date{}

\begin{document}

\hypertarget{ux89e3ux7b54-01-ux30d9ux30afux30c8ux30ebux3068ux30d9ux30afux30c8ux30ebux7a7aux9593}{%
\section{解答 01 ---
ベクトルとベクトル空間}\label{ux89e3ux7b54-01-ux30d9ux30afux30c8ux30ebux3068ux30d9ux30afux30c8ux30ebux7a7aux9593}}

\hypertarget{section}{%
\subsection{1}\label{section}}

\texttt{av\_1+bv\_2=(a+b,a,b)\^{}T} です。\texttt{(3,1,2)\^{}T}
と比較すると \texttt{a=1,b=2,a+b=3} が同時に成立するので属します。

重要なのは、これは \texttt{Ax=b} を解いていることと同じだという点です。

\hypertarget{section-1}{%
\subsection{2}\label{section-1}}

\texttt{v\_2=2v\_1} なので

\[
-2v_1+v_2+0v_3=0
\]

という非自明な一次結合が存在します。したがって一次従属です。

\hypertarget{section-2}{%
\subsection{3}\label{section-2}}

\texttt{R\^{}3} の次元は3です。3本の一次独立ベクトルの span
は3次元部分空間になります。\texttt{R\^{}3} の3次元部分空間は
\texttt{R\^{}3} 自身なので、空間全体を張ります。よって基底です。

\hypertarget{section-3}{%
\subsection{4}\label{section-3}}

\texttt{z=-x-y} より

\[
(x,y,z)^T=x(1,0,-1)^T+y(0,1,-1)^T.
\]

したがって基底の一例は

\[
\{(1,0,-1)^T,(0,1,-1)^T\}
\]

で、次元は2です。零ベクトルを含み、加法とスカラー倍で条件
\texttt{x+y+z=0} が保存されるため部分空間です。

\hypertarget{section-4}{%
\subsection{5}\label{section-4}}

ベクトルは抽象的・幾何学的対象で、座標は基底に依存する数値表現です。基底を変えると座標成分は変わりますが、ベクトルそのもの、一次独立性、空間の次元などは変わりません。この視点が対角化や
SVD の「良い基底を選ぶ」という発想につながります。

\end{document}
